\documentclass[12pt]{article}
\usepackage[utf8]{inputenc}
\usepackage{amsmath,amssymb}
\usepackage{geometry}
\usepackage{setspace}
\usepackage{enumitem}
\geometry{margin=1in}
\setstretch{1.15}

\begin{document}

\begin{center}
    {\LARGE \textbf{Lecture Scribe (Exam Revision Format)}}\\[1em]
    {\Large \textbf{CSE400 – Fundamentals of Probability in Computing}}\\[0.5em]
    {\large \textbf{Lecture L15\_16\_17\_18: Joint Random Variables + Joint Distributions + Examples}}
\end{center}

\vspace{0.5em}

\noindent \textbf{Source used:} Lecture PDF only

\hrule
\vspace{1em}

\section*{1. Introduction: Why Study Joint Random Variables?}

In science and in real life, we are often interested in \textbf{two (or more) random variables at the same time}. The lecture motivates this using several examples:

\begin{itemize}
    \item Height and weight of giraffes
    \item IQ and birthweight of children
    \item Frequency of exercise and rate of heart disease
    \item Air pollution level and respiratory illness rate
    \item Number of Facebook friends and age of members
\end{itemize}

The lecture also provides application-driven examples:

\subsection*{Example 1: Smart Transportation}

\begin{itemize}
    \item \(X =\) Vehicle Speed
    \item \(Y =\) Traffic Density
\end{itemize}

\noindent \textbf{Question posed in lecture:}

Can speed be high when the traffic density is also high?

\noindent \textbf{Answer in lecture:} Yes / No (posed as a conceptual relationship question)

This motivates that \(X\) and \(Y\) may be \textbf{correlated random variables}. Applications mentioned:

\begin{itemize}
    \item Adaptive traffic signals
    \item Congestion prediction
    \item Autonomous vehicle routing
\end{itemize}

\subsection*{Example 2: Wireless Network Performance}

\begin{itemize}
    \item \(X =\) Signal-to-Noise Ratio (SNR)
    \item \(Y =\) Data Throughput
\end{itemize}

\noindent \textbf{Reason given in lecture:}

Throughput depends on SNR, but both vary due to fading, interference, and congestion.

\noindent \textbf{Question posed in lecture:}

If SNR is high, will throughput always be high?

\noindent \textbf{Answer in lecture:} Yes / No

\noindent \textbf{Reasoning explicitly stated in lecture:}

Network congestion (conditional relationships).

Applications mentioned:

\begin{itemize}
    \item 5G scheduling
    \item Adaptive modulation
    \item Network planning
\end{itemize}

\subsection*{Example 3: Weather Prediction}

\begin{itemize}
    \item \(X =\) Rainfall
    \item \(Y =\) Temperature
\end{itemize}

Application: weather forecasting.

\subsection*{Example 4: Drone Delivery}

\begin{itemize}
    \item \(X =\) Wind Speed
    \item \(Y =\) Delivery Time
\end{itemize}

Application: logistics and UAV systems.

\subsection*{Example 5: Rolling Two Dice}

The lecture specifically mentions events such as:

\begin{itemize}
    \item \(X+Y=7\)
    \item \(X>Y\)
\end{itemize}

\noindent \textbf{Conclusion stated in lecture:}

This requires a \textbf{joint PMF}.

\hrule
\vspace{1em}

\subsection*{Key Conceptual Conclusion}

The lecture states:

\begin{itemize}
    \item In such situations, the random variables have a \textbf{joint distribution}.
    \item A joint distribution allows us to:
    \begin{enumerate}
        \item Compute probabilities of events involving both variables.
        \item Understand the relationship between the variables.
    \end{enumerate}
    \item This is simplest when the variables are \textbf{independent}.
    \item When they are not independent, we use \textbf{covariance} and \textbf{correlation} as measures of the nature of dependence between them.
\end{itemize}

\hrule
\vspace{1em}

\section*{2. Discrete Case – Joint PMF}

\subsection*{2.1 Setup}

Suppose \(X\) and \(Y\) are discrete random variables.

The lecture gives:

\begin{itemize}
    \item \(X \in \{x_1, x_2, \dots, x_n\}\)
    \item \(Y \in \{y_1, y_2, \dots, y_m\}\)
\end{itemize}

Then the ordered pair \((X,Y)\) takes values in the \textbf{product set}:

\[
\{(x_1,y_1),(x_1,y_2),\dots,(x_n,y_m)\}
\]

This means every possible joint outcome is a pair formed by one value of \(X\) and one value of \(Y\).

\hrule
\vspace{1em}

\subsection*{2.2 Definition of Joint PMF}

\subsubsection*{Definition (Joint PMF)}

The lecture defines the joint PMF as:

\[
p(x_i,y_j)=P(X=x_i,\;Y=y_j)
\]

\subsubsection*{Meaning}

This gives the probability that the \textbf{joint outcome occurs simultaneously}, i.e., both:

\begin{itemize}
    \item \(X=x_i\), and
    \item \(Y=y_j\).
\end{itemize}

\hrule
\vspace{1em}

\section*{3. Joint PMF Table Representation}

The lecture states that the joint PMF is organized in a \textbf{table}, where:

\begin{itemize}
    \item rows correspond to values of one variable,
    \item columns correspond to values of the other variable,
    \item each cell contains the joint probability \(p(x_i,y_j)\).
\end{itemize}

\hrule
\vspace{1em}

\subsection*{3.1 Key Properties of the Joint PMF}

The lecture explicitly lists the following properties:

\subsubsection*{Property 1: Bounds}

\[
0 \le p(x,y) \le 1
\]

\subsubsection*{Property 2: Total Probability}

\[
\sum_i \sum_j p(x_i,y_j)=1
\]

This means the probabilities across \textbf{all possible joint outcomes} must sum to 1.

\hrule
\vspace{1em}

\section*{4. Example: Rolling Two Dice (Discrete Joint PMF Example)}

\subsection*{4.1 Experiment Setup}

The lecture defines the experiment:

\begin{itemize}
    \item Roll two fair dice.
\end{itemize}

Define:

\begin{itemize}
    \item \(X =\) value on the \textbf{first die}
    \item \(T =\) \textbf{total (sum)} of both dice
\end{itemize}

\hrule
\vspace{1em}

\subsection*{4.2 Probability Structure}

The lecture states:

\begin{itemize}
    \item Each outcome \((i,j)\) from the two dice has probability:
\end{itemize}

\[
\frac{1}{36}
\]

Therefore, the joint probability table of \((X,T)\) can be constructed from the equally likely two-die outcomes.

\hrule
\vspace{1em}

\subsection*{4.3 Observations from the Lecture}

The lecture explicitly states:

\begin{enumerate}
    \item \textbf{Entries are either \(0\) or \(\frac{1}{36}\)}
    \item \textbf{\(X\) and \(T\) are not independent}
\end{enumerate}

\hrule
\vspace{1em}

\subsection*{4.4 Exam-Style Interpretation of the Table (Exactly as implied by lecture)}

\subsubsection*{Step 1: Identify the variables}

\begin{itemize}
    \item \(X\) is the first die value.
    \item \(T\) is the sum of the two dice.
\end{itemize}

\subsubsection*{Step 2: Determine the sample space}

Each ordered pair \((i,j)\), where:

\begin{itemize}
    \item \(i \in \{1,2,3,4,5,6\}\)
    \item \(j \in \{1,2,3,4,5,6\}\)
\end{itemize}

has probability \(\frac{1}{36}\).

\subsubsection*{Step 3: Translate to \((X,T)\)}

For a fixed first die value \(X=i\), the total is:

\[
T=i+j
\]

where \(j\) is the second die value.

\subsubsection*{Step 4: Determine valid joint outcomes}

A joint pair \((X,T)=(i,t)\) is possible only if:

\[
t=i+j \quad \text{for some } j \in \{1,2,3,4,5,6\}
\]

Equivalently,

\[
j=t-i \in \{1,2,3,4,5,6\}
\]

\subsubsection*{Step 5: Assign probability}

\begin{itemize}
    \item If \(t-i \in \{1,2,3,4,5,6\}\), then exactly one second-die value produces that pair, so:
\end{itemize}

\[
P(X=i,T=t)=\frac{1}{36}
\]

\begin{itemize}
    \item Otherwise:
\end{itemize}

\[
P(X=i,T=t)=0
\]

\subsubsection*{Step 6: Conclusion}

Hence the lecture’s observation follows:

\begin{itemize}
    \item all table entries are either \(0\) or \(\frac{1}{36}\),
    \item and \(X\) and \(T\) are not independent.
\end{itemize}

\hrule
\vspace{1em}

\section*{5. Continuous Case – Joint PDF}

\subsection*{5.1 Analogy with the Discrete Case}

The lecture explicitly states the following analogies:

\begin{itemize}
    \item \textbf{Discrete values} \(\rightarrow\) \textbf{Continuous intervals}
    \item \textbf{Joint PMF} \(\rightarrow\) \textbf{Joint PDF}
    \item \textbf{Sums} \(\rightarrow\) \textbf{Double integrals}
\end{itemize}

\hrule
\vspace{1em}

\subsection*{5.2 Setup}

If:

\begin{itemize}
    \item \(X \in [a,b]\)
    \item \(Y \in [c,d]\)
\end{itemize}

then the pair \((X,Y)\) takes values in the \textbf{product region}:

\[
[a,b]\times[c,d]
\]

This is the rectangular region formed by the ranges of \(X\) and \(Y\).

\hrule
\vspace{1em}

\subsection*{5.3 Definition of Joint PDF}

\subsubsection*{Definition (Joint PDF)}

A function \(f(x,y)\) is called the \textbf{joint probability density function} of \(X\) and \(Y\) if:

\[
P((X,Y)\text{ lies in a small rectangle of area }dx\,dy)=f(x,y)\,dx\,dy
\]

This is the lecture’s formal density interpretation.

\hrule
\vspace{1em}

\subsection*{5.4 Key Properties of the Joint PDF}

The lecture explicitly lists:

\subsubsection*{Property 1: Nonnegativity}

\[
f(x,y)\ge 0
\]

\subsubsection*{Property 2: Normalization}

\[
\iint f(x,y)\,dx\,dy = 1
\]

The lecture also emphasizes:

\begin{itemize}
    \item The joint PDF \(f(x,y)\) \textbf{may take values greater than 1}.
    \item It is a \textbf{probability density}, \textbf{not a probability}.
\end{itemize}

\hrule
\vspace{1em}

\section*{6. Worked Example (Continuous Joint PDF)}

\subsection*{6.1 Given Data}

The lecture gives:

\[
f(x,y)=4xy,\qquad 0<x<1,\;0<y<1
\]

The lecture states that jointly, \((X,Y)\) takes values in the \textbf{unit square}:

\[
[0,1]^2
\]

\subsubsection*{Tasks in the lecture}

\begin{enumerate}
    \item Show that \(f(x,y)\) is a valid joint PDF.
    \item Visualize the event
    \[
    A=\{X<0.5,\;Y>0.5\}
    \]
    \item Compute \(P(A)\).
\end{enumerate}

\hrule
\vspace{1em}

\section*{7. Example Solution – Step-by-Step (Exactly Following Lecture)}

\subsection*{7.1 Step 1: Check Validity (Normalization)}

The lecture computes:

\[
\int_0^1 \int_0^1 4xy\,dx\,dy
\]

\subsubsection*{Step-by-step derivation exactly in exam style}

\subsubsection*{Step 1.1: Write the normalization integral}

\[
\int_0^1 \int_0^1 4xy\,dx\,dy
\]

\subsubsection*{Step 1.2: Integrate with respect to \(x\) first}

\[
= \int_0^1 \left( \int_0^1 4xy\,dx \right) dy
\]

Since \(y\) is constant with respect to \(x\):

\[
\int_0^1 4xy\,dx = 4y\int_0^1 x\,dx
\]

\[
= 4y\left[\frac{x^2}{2}\right]_0^1
\]

\[
= 4y\left(\frac{1}{2}\right)=2y
\]

So:

\[
\int_0^1 \int_0^1 4xy\,dx\,dy = \int_0^1 2y\,dy
\]

\subsubsection*{Step 1.3: Integrate with respect to \(y\)}

\[
\int_0^1 2y\,dy = [y^2]_0^1
\]

\[
=1^2-0^2=1
\]

\subsubsection*{Conclusion from the lecture}

Since:

\begin{itemize}
    \item \(f(x,y)>0\), and
    \item total probability is 1,
\end{itemize}

therefore:

\[
f(x,y)=4xy
\]

is a \textbf{valid joint PDF}.

\hrule
\vspace{1em}

\subsection*{7.2 Step 2: Compute \(P(A)\)}

The event is:

\[
A=\{X<0.5,\;Y>0.5\}
\]

So the lecture computes probability over:

\begin{itemize}
    \item \(0<x<0.5\)
    \item \(0.5<y<1\)
\end{itemize}

\subsubsection*{Step-by-step derivation exactly as presented}

\subsubsection*{Step 2.1: Write the probability integral}

\[
P(A)=\int_0^{0.5}\int_{0.5}^{1}4xy\,dy\,dx
\]

\subsubsection*{Step 2.2: Integrate with respect to \(y\) first}

\[
= \int_0^{0.5}\left[2xy^2\right]_{0.5}^{1}dx
\]

\subsubsection*{Step 2.3: Substitute the bounds}

\[
= \int_0^{0.5}2x(1-0.25)\,dx
\]

\[
= \int_0^{0.5}2x(0.75)\,dx
\]

\[
= \int_0^{0.5}\frac{3}{2}x\,dx
\]

\subsubsection*{Step 2.4: Integrate with respect to \(x\)}

\[
= \left[\frac{3}{4}x^2\right]_0^{0.5}
\]

\[
= \frac{3}{4}(0.5)^2
\]

\[
= \frac{3}{4}\cdot \frac{1}{4}
\]

\[
= \frac{3}{16}
\]

\subsubsection*{Final Answer}

\[
P(A)=\frac{3}{16}
\]

This is the exact result shown in the lecture.

\hrule
\vspace{1em}

\section*{8. Joint Cumulative Distribution Function (Joint CDF)}

\subsection*{8.1 Notation}

Suppose \(X\) and \(Y\) are jointly distributed random variables.

The lecture states the notation:

\[
X<x,\;Y<y
\]

means the event:

\[
\{X<x \text{ and } Y<y\}
\]

This is the event that \textbf{both inequalities occur simultaneously}.

\hrule
\vspace{1em}

\subsection*{8.2 Definition of Joint CDF}

\subsubsection*{Definition}

The joint cumulative distribution function is defined as:

\[
F(x,y)=P(X<x,\;Y<y)
\]

This is the probability that:

\begin{itemize}
    \item \(X\) is less than \(x\), and
    \item \(Y\) is less than \(y\).
\end{itemize}

\hrule
\vspace{1em}

\section*{9. Joint CDF – Continuous Case}

If \(X\) and \(Y\) are continuous random variables with joint density \(f(x,y)\) over \([a,b]\times[c,d]\), then the lecture states that the joint CDF is obtained by \textbf{integrating the density}:

\[
F(x,y)=\int_c^y\int_a^x f(u,v)\,du\,dv
\]

This expresses cumulative probability over the rectangle from the lower-left corner of the support up to \((x,y)\).

\hrule
\vspace{1em}

\subsection*{9.1 Recovering the Joint PDF from the Joint CDF}

The lecture explicitly states:

Since there are two variables, we use \textbf{partial derivatives}:

\[
f(x,y)=\frac{\partial^2}{\partial x\,\partial y}F(x,y)
\]

This is the recovery formula from the continuous joint CDF.

\hrule
\vspace{1em}

\section*{10. Joint CDF – Discrete Case}

If \(X\) and \(Y\) are discrete random variables with joint PMF \(p(x_i,y_j)\), then the lecture states that the joint CDF is obtained by \textbf{summing probabilities}:

\[
F(x,y)=\sum_{x_i<x}\sum_{y_j<y} p(x_i,y_j)
\]

This means we add all joint PMF entries whose:

\begin{itemize}
    \item row values satisfy \(x_i<x\), and
    \item column values satisfy \(y_j<y\).
\end{itemize}

\hrule
\vspace{1em}

\subsection*{10.1 Interpretation (Exactly as Given in Lecture)}

The lecture gives the following interpretation:

\begin{itemize}
    \item Add all probabilities in the table
    \item For rows with \(x_i<x\)
    \item And columns with \(y_j<y\)
\end{itemize}

\hrule
\vspace{1em}

\section*{11. Properties of the Joint CDF}

The lecture states that the joint CDF \(F(x,y)\) of \(X\) and \(Y\) must satisfy several properties.

\hrule
\vspace{1em}

\subsection*{11.1 Property 1: Non-decreasing}

\[
F(x,y) \text{ is non-decreasing}
\]

\subsubsection*{Meaning stated in lecture}

If \(x\) or \(y\) increase, then \(F(x,y)\) must:

\begin{itemize}
    \item stay constant, or
    \item increase.
\end{itemize}

\hrule
\vspace{1em}

\subsection*{11.2 Property 2: Lower-left Limit}

The lecture states:

\[
F(x,y)=0 \text{ at the lower-left of the joint range}
\]

If the lower-left corner is \((-\infty,-\infty)\), then:

\[
\lim_{(x,y)\to(-\infty,-\infty)} F(x,y)=0
\]

This expresses that cumulative probability before the support begins is zero.

\hrule
\vspace{1em}

\subsection*{11.3 Property 3: Upper-right Limit}

The lecture states:

\[
F(x,y)=1 \text{ at the upper-right of the joint range}
\]

If the upper-right corner is \((\infty,\infty)\), then:

\[
\lim_{(x,y)\to(\infty,\infty)} F(x,y)=1
\]

This expresses that cumulative probability over the entire joint range equals 1.

\hrule
\vspace{1em}

\section*{12. Consolidated Exam Revision Summary}

\subsection*{12.1 Core Definitions to Memorize}

\subsubsection*{Discrete Joint PMF}

\[
p(x_i,y_j)=P(X=x_i,\;Y=y_j)
\]

\subsubsection*{Continuous Joint PDF}

\[
P((X,Y)\text{ lies in a small rectangle of area }dx\,dy)=f(x,y)\,dx\,dy
\]

\subsubsection*{Joint CDF}

\[
F(x,y)=P(X<x,\;Y<y)
\]

\hrule
\vspace{1em}

\subsection*{12.2 Key Discrete Facts}

\begin{itemize}
    \item \((X,Y)\) takes values in the \textbf{product set} of values of \(X\) and \(Y\).
    \item Joint PMF table entries represent simultaneous outcomes.
    \item Properties:
\end{itemize}

\[
0\le p(x,y)\le 1
\]

\[
\sum_i\sum_j p(x_i,y_j)=1
\]

\hrule
\vspace{1em}

\subsection*{12.3 Key Continuous Facts}

\begin{itemize}
    \item \((X,Y)\) takes values in a \textbf{product region} such as \([a,b]\times[c,d]\).
    \item Properties:
\end{itemize}

\[
f(x,y)\ge 0
\]

\[
\iint f(x,y)\,dx\,dy=1
\]

\begin{itemize}
    \item \(f(x,y)\) may be greater than 1 because it is a \textbf{density}, not a probability.
\end{itemize}

\hrule
\vspace{1em}

\subsection*{12.4 Joint CDF Formulas}

\subsubsection*{Continuous}

\[
F(x,y)=\int_c^y\int_a^x f(u,v)\,du\,dv
\]

\subsubsection*{Recover density}

\[
f(x,y)=\frac{\partial^2}{\partial x\,\partial y}F(x,y)
\]

\subsubsection*{Discrete}

\[
F(x,y)=\sum_{x_i<x}\sum_{y_j<y}p(x_i,y_j)
\]

\hrule
\vspace{1em}

\subsection*{12.5 Worked Example Final Results}

For:

\[
f(x,y)=4xy,\quad 0<x<1,\;0<y<1
\]

\begin{itemize}
    \item Validity check:
\end{itemize}

\[
\int_0^1\int_0^1 4xy\,dx\,dy=1
\]

Hence valid joint PDF.

\begin{itemize}
    \item For:
\end{itemize}

\[
A=\{X<0.5,\;Y>0.5\}
\]

\[
P(A)=\int_0^{0.5}\int_{0.5}^{1}4xy\,dy\,dx=\frac{3}{16}
\]

\hrule
\vspace{1em}

\section*{13. Final Exam-Oriented Closing Note}

This lecture establishes the foundational framework for analyzing \textbf{two random variables simultaneously}:

\begin{enumerate}
    \item \textbf{Discrete case:} Joint PMF and table representation
    \item \textbf{Continuous case:} Joint PDF and double integrals
    \item \textbf{Cumulative form:} Joint CDF in both continuous and discrete settings
    \item \textbf{Worked examples:}
    \begin{itemize}
        \item Two-dice joint PMF structure
        \item Continuous joint PDF example with full normalization and event probability computation
    \end{itemize}
\end{enumerate}

All formulas, interpretations, and derivations above are reproduced strictly from the provided lecture material and organized in step-by-step exam style for revision.

\hrule

\end{document}